\itemhead{习题2.2 · 第5题（教材p.49）}{XJ-B1-C02-HW-P049-Y0879}
\label{XJ-B1-C02-HW-P049-Y0879}
{\footnotesize 湘教版教材p.49 / PDF第56页 / 习题2.2 / 第5题}
\field{原题}{已知二次函数 $y=(1-m)x^2+4x-3$ 的图象开口向下，与 $x$ 轴交于 $A(x_1,0),B(x_2,0)$ 两点。
\subq{1}求 $m$ 的取值范围；
\subq{2}当 $x_1^2+x_2^2=10$ 时，求该二次函数的表达式。}
\field{参考答案}{（1）$1<m<7/3$；（2）由根和、根积得 $16/(1-m)^2+6/(1-m)=10$，结合（1）取 $m=2$，故 $y=-x^2+4x-3$。}
{\small \field{关联知识}{Q01 二次函数、零点、判别式与根系关系（教材帮p.65 / PDF74）}
\field{方法与易错点}{按已知零点、顶点或系数选择交点式、顶点式或一般式；根平方和用根系关系消元。 易错：遗漏 $a
e0$；把二重根算成两个不同零点；有两正根只查积正。}
\field{答案依据}{答案册相关题组：PDF第13、14页（册内页12、13）；按当前题库原题逐项重算，不按旧版题号机械对套。}}
